% Part: sets-functions-relations
% Chapter: arithmetization
% Section: cauchy

\documentclass[../../../include/open-logic-section]{subfiles}

\begin{document}

\olfileid{sfr}{arith}{cauchy}

\olsection{附录：实数作为柯西序列}

在\olref[cuts]{sec}中，我们用戴德金分割构造了实数。
本节将解释另一种构造方法。它基于 Cauchy（柯西）对（即我们现在所称的）柯西序列的定义；
但用这一定义来\emph{构造}实数则要归功于其他十九世纪的学者，特别是 Weierstrass、Heine（海涅）、M\'{e}ray（梅雷）和 Cantor。（一段精彩的历史介绍，可参阅 \citeauthor{OConnorRobertson:RN} \citeyear{OConnorRobertson:RN}。）

在进入十九世纪之前，值得先提一下 Simon Stevin (1548--1620)。简而言之，Stevin 认识到我们可以用十进制展开来理解每个实数。因此，即使是像 $\sqrt{2}$ 这样的无理数，也有一个从如下开始的十进制展开：
\[
	1.41421356237\ldots
\] 
在集合论中表示十进制展开非常容易：只需把它们看作函数 $d \colon \Nat \to \Nat$，其中 $d(n)$ 是我们感兴趣的第 $n$ 位小数。我们还需要稍作调整，以处理实数中小数点之前的部分（此处就是 $1$），并需引入一个等价关系来保证，例如，$0.999\ldots = 1$ 成立。但以 Stevin 的方式，在集合论内部给出一个完全严格的实数构造并不困难。

Stevin 并非我们的重点。（更多关于 Stevin 的信息，参见 \citealt{KatzKatz2012}。）但这里有一个密切相关的想法。不直接处理 $\sqrt{2}$ 的十进制展开，我们可以转而考虑一个对 $\sqrt{2}$ 越来越精确的有理数逼近序列，即考虑越来越精确的展开：
\[
1, 1.4, 1.414, 1.4142, 1.41421,\ldots
\]
通过“越来越好的逼近”来考虑实数的想法，为我们提供了另一系列洞见的基础（类似于我们从 $\Int$ 构造 $\Rat$，或从 $\Nat$ 构造 $\Int$ 时的那些领悟）。基本的洞见如下：
\begin{enumerate}
	\item 每个实数都可以写成一个（可能是无限的）十进制展开。
	\item 一个（可能是无限的）十进制展开所编码的信息，同样可以由一个有理数序列来编码。
	\item 一个有理数序列可以看作一个从 $\Nat$ 到 $\Rat$ 的函数；只需令 $f(n)$ 为序列中的第 $n$ 个有理数。
\end{enumerate}
当然，并非\emph{任意}一个从 $\Nat$ 到 $\Rat$ 的函数都能给我们一个实数。例如，考虑这个函数：
\[
	f(n) =\begin{cases}
			1 & \text{如果 }n\text{ 是奇数}\\
			0 &\text{如果 }n\text{ 是偶数}
		\end{cases}
\]
本质上，问题在于序列 $0,1,0,1,0,1,0,\ldots$
似乎并不“趋近于”任何实数。因此，为了确保我们考虑的序列确实趋近于某个实数，我们需要把注意力限制在具有某种\emph{极限}的序列上。

我们已经在\olref[his][set][limits]{sec}中遇到过极限的概念。但我们不能使用\emph{完全}相同的定义。那里的表达式“$(\forall \epsilon>0)$”隐含了对实数的量化；而且我们当时考虑的是实数上函数的极限；所以援引那个定义就等于预设了实数的存在；而实数正是我们想要\emph{构造}的对象。幸运的是，我们可以使用一个密切相关的极限概念。
\begin{defn}\ollabel{def:CauchySequence}
  一个函数 $f: \Nat \to \Rat$ 是\emph{柯西序列}，当且仅当对于任意正有理数 $\epsilon \in \Rat$，都有 $(\exists \ell \in
  \Nat)(\forall m, n > \ell)|f(m) - f(n)| < \epsilon$。
\end{defn}

极限的基本思想与之前相同：如果你想要某个精确度（用~$\epsilon$ 衡量），就存在一个需要查看的“区域”（任何大于~$\ell$ 的输入）。并且很容易看出我们的序列 $1$, $1.4$, $1.414$, $1.4142$,
$1.41421$\ldots 有一个极限：如果你想在 $\nicefrac{1}{10^{n}}$ 的误差内逼近 $\sqrt{2}$，
那么只需查看第 $n$ 项之后的任意一项。

那么，一个自然的想法是，实数就\emph{是}某个柯西序列。但是，正如构造 $\Int$ 和 $\Rat$ 时那样，这会过于天真：对于任意给定的实数，都有多个不同的柯西序列指示该实数。一个简单的方法可以说明这一点。给定一个柯西序列~$f$，定义 $g$ 为一个除 $g(0)\neq
f(0)$ 外与~$f$ 完全相同的函数。由于这两个序列在第一个数之后处处相同，我们（最终）会想说它们有相同的极限（就 \olref{def:CauchySequence} 中采用的极限意义而言），因此应视为“定义”了同一个实数。所以，我们确实应该把这些柯西序列看作同一个实数。

因此，我们再次需要在柯西序列上定义一个等价关系，并将实数等同于等价类。首先我们需要函数趋于 $0$ 的概念。对于任意函数 $h : \Nat \to \Rat$，称\emph{$h$ 趋于 $0$}当且仅当对任意正的有理数 $\epsilon \in \Rat$，我们有 $(\exists \ell \in \Nat)(\forall n > \ell)|h(n)| <
\epsilon$。\footnote{将此定义与 \olref[his][set][limits]{sec} 中 $\lim_{x
\mathord{\rightarrow}\infty}f(x) = 0$ 的定义进行比较。} 进一步地，当 $f$ 和 $g$ 是函数 $\Nat \to \Rat$ 时，令 $(f-g)(n) = f(n) - g(n)$。现在定义：
\[
	f \Realequiv g \text{ 当且仅当 $(f-g)$ 趋于 $0$}.
\]
我们需要检查 $\Realequiv$ 是一个等价关系；确实如此。然后，我们不妨将实数定义为所有从 $\Nat$ 到
$\Rat$ 的柯西序列在 $\Realequiv$ 下的等价类。

\begin{prob}
设对于每个 $n$，$f(n) = 0$。令 $g(n) = \frac{1}{(n+1)^2}$。证明
两者都是柯西序列，并且实际上两个函数的极限都是 $0$，从而也有 $f \sim_\Real g$。
\end{prob}

如同通常做法，我们将包含 $f$ 的等价类记作 $\equivrep{f}{\Realequiv}$。然而，为了保持可读性，在下文中我们将省略下标而只写作 $\equivrep{f}{}$。我们还约定，对于每个 $q \in \Rat$，$q_{\Real} = \equivrep{c_{q}}{}$，其中 $c_{q}$ 是对所有 $n \in \Nat$ 满足 $c_q(n) = q$ 的常值函数。然后我们定义实数上的基本关系和运算，例如：
\begin{align*}
	\equivrep{f}{} + \equivrep{g}{} &= 	\equivrep{(f + g)}{} \\
	\equivrep{f}{} \times 	\equivrep{g}{} &= \equivrep{(f \times g)}{} 
\end{align*}
其中 $(f + g)(n) = f(n) + g(n)$ 且 $(f \times g)(n) = f(n) \times
g(n)$。当然，我们还需要检查当 $f$ 和 $g$ 是柯西序列时， $(f + g)$、
$(f-g)$ 和 $(f\times g)$ 中的每一个也都是柯西序列；确实如此，我们将这留给你来完成。

最后，我们定义序的概念。称 $\equivrep{f}{}$ 是\emph{正的}，如果 $\equivrep{f}{}\neq 0_\Rat$ 且 $(\exists
\ell \in \Nat)(\forall n > \ell)0 < f(n)$。然后称 $\equivrep{f}{} <
\equivrep{g}{}$ 当且仅当 $\equivrep{(g - f)}{}$ 是正的。我们必须检查这是良定义的（即，它不依赖于等价类中“代表元”的选择）。
%\begin{align*}
%	\equivrep{f}{} \leq \equivrep{g}{} \text{ iff }&\equivrep{f}{} = \equivrep{g}{} \text{ or }\\
%	&(\exists \ell \in \Nat)(\forall n \in \Nat)(\ell < n \lif f(n) \leq g(n))
%\end{align*}
完成验证后，很容易证明这些运算和序给出了预期的代数性质；即：
\begin{thm}\ollabel{thm:cauchyorderedfield}
柯西序列构成一个有序域。
\end{thm}

\begin{proof}
练习。
\end{proof}

\begin{prob}
证明柯西序列构成一个有序域。
\end{prob}

更困难的是证明如此构造的实数具有完备性，因此我们将给出证明。

\begin{thm}
每个有上界的非空柯西序列集都有最小上界。
\end{thm}

\begin{proof}[证明思路]
设 $S$ 为任意有上界的非空柯西序列集。那么存在某个 $p \in \Rat$，使得 $p_{\Real}$ 是 $S$ 的一个上界。令 $r \in S$；则存在某个 $q \in \Rat$ 使得 $q_{\Real} < r$。因此，如果 $S$ 存在最小上界，那么它介于 $q_\Real$ 和 $p_\Real$ 之间（含端点）。

我们将通过从下方和上方同时逼近来锁定这个最小上界。具体来说，我们定义两个函数 $f, g \colon
\Nat \to \Rat$，目标是让 $f$ 从上方逼近最小上界，而 $g$ 从下方逼近它。我们首先定义：
\begin{align*}
	f(0) &= p \\
	g(0) &= q
\end{align*}
然后，当 $a_n = \frac{f(n) + g(n)}{2}$ 时，令：\footnote{这是一个递归定义。但我们\emph{尚未}给出任何理由来认为递归定义是合法的。}
\begin{align*}
	f(n+1) &=
	\begin{cases}
		a_n &\text{若 }(\forall h \in S)\equivrep{h}{} \leq (a_n)_\Real\\
		f(n)&\text{否则}
	\end{cases}\\
	g(n+1) &=
	\begin{cases}
		a_n &\text{若 }(\exists h \in S)\equivrep{h}{} \geq (a_n)_\Real\\
	 	g(n) &\text{否则}
	\end{cases}
\end{align*}
$f$ 和 $g$ 都是柯西序列。（这可以很容易地验证，但我们留作练习。）注意函数 $(f-g)$ 趋于 $0$，因为 $f$ 和 $g$ 之间的差值每一步都减半。因此 $\equivrep{f}{} = \equivrep{g}{}$。

我们首先证明 $\equivrep{f}{}$ 是 $S$ 的一个上界，即 $(\forall h \in S)\equivrep{h}{} \leq \equivrep{f}{}$。（我们将在过程中引用 \olref{thm:cauchyorderedfield}。）令 $h \in S$ 并假设（为归谬）$\equivrep{f}{} < \equivrep{h}{}$，从而 $0_\Real < \equivrep{(h-f)}{}$。由于 $f$ 是一个单调递减的柯西序列，存在某个 $n \in \Nat$ 使得 $\equivrep{(c_{f(n)} - f)}{} < \equivrep{(h-f)}{}$。所以：
\[
	(f(n))_\Real = \equivrep{c_{f(n)}}{} < \equivrep{f}{} + \equivrep{(h-f)}{} = \equivrep{h}{},
\]
这与以下事实矛盾：根据构造，$\equivrep{h}{} \leq (f(n))_\Real$。

接下来我们证明 $\equivrep{f}{} = \equivrep{g}{}$ 是 $S$ 的\emph{最小}上界。令 $j$ 为任意柯西序列并假设 $\equivrep{j}{} < \equivrep{g}{}$。与上面的推理类似（利用 $g$ 是\emph{递增}的事实），存在 $n \in \Nat$ 使得 $\equivrep{j}{} < (g(n))_\Real$。但由构造，存在 $h \in S$ 使得 $(g(n))_\Real \leq \equivrep{h}{}$，所以 $\equivrep{j}{} < \equivrep{h}{}$，因此 $\equivrep{j}{}$ 不是 $S$ 的上界。
\end{proof}

\end{document}